%% Spaltendefinition:
%% Erste Spalte vom Typ "l": Spalte so breit, wie der Inhalt, linksbündige Ausrichtung
%% Drei weitere Spalten vom Typ "y": die Breite wird automatisch berechnet
%% in Abhängigkeit von der Gesamtbreite der Tabelle und Breite anderer Spalten,
%% Inhalt wird umbrochen, wenn zu lang
%% Format: rechtsbündiger Flattersatz ohne Silbentrennung, 
%% Weitere Möglichkeiten:
%% "X" = Blocksatz (mit Silbentrennung)
%% "x" = linksbündiger Flattersatz ohne Silbentrennung
%% "F" = linksbündiger Flattersatz mit Silbentrennung (KSP-Verlag: unerwünscht),
%% "Y" = rechtsbündiger Flattersatz mit Silbertrennung
\begin{longtable}{l*{3}{y}}%
%% Definition für die erste Kopfzeile
	%% Bei Longtable ist Caption Teil der Kopfzeile
	\hiderowcolors%
	%\vspace{-1.6\baselineskip}% Workaround, um den Abstand vor der Überschrift, welcher durch arraystretch eingeführt wird, zu "entfernen"
	%\vspace{-18pt}%
	%\vspace{\dimexpr (\baselineskip * \numexpr (1 - \numexpr \arraystretch \relax) \relax) \relax}}%
	\tabularnewline%
	\caption{Longtable Tabelle über einen Seitenumbruch hinaus.}%
	\label{tab:MehrseitigeTabelle}%
	%\vspace{\mytablecaptionskip}%
	\tabularnewline%
	%\\%
	%% Definition für die erste Kopfzeile
	\rowcolor{tableheadcolor}%
	\toprule%
	\tableheadfont Tabellenkopf &
	\tableheadfont Tabellenkopf	&
	\tableheadfont Tabellenkopf	&
	\tableheadfont Tabellenkopf \tabularnewline
	\midrule%
\endfirsthead
%% Definition für die Kopfzeile auf den Folgeseiten
	\rowcolor{tableheadcolor}%
	\toprule%
	\tableheadfont Zwischenkopf &
	\tableheadfont Zwischenkopf	&
	\tableheadfont Zwischenkopf	&
	\tableheadfont Zwischenkopf \tabularnewline
	\midrule%
\endhead
%% Definition für die Fußzeile auf Zwischenseiten
	\hiderowcolors%
	\bottomrule%
	\multicolumn{4}{r}{\small\emph{weiter auf der nächsten Seite}}
\endfoot
%% Definition für das Tabellenende
	%\showrowcolors%
	\bottomrule%
\endlastfoot
%\tablestylealtcolored%
\showrowcolors%
\bfseries Beschreibung 	& 1 Inhalt & Inhalt & Inhalt \tabularnewline
\bfseries Beschreibung   & 2 Inhalt & Inhalt & Inhalt \tabularnewline
\bfseries Beschreibung   & 3 Inhalt & Inhalt & Inhalt \tabularnewline
\bfseries Beschreibung   & 4 Inhalt & Inhalt & Inhalt \tabularnewline
\bfseries Beschreibung   & 5 Inhalt & Inhalt & Inhalt \tabularnewline
\bfseries Beschreibung   & 6 Inhalt & Inhalt & Inhalt \tabularnewline
\bfseries Beschreibung   & 7 Inhalt & Inhalt & Inhalt \tabularnewline
\bfseries Beschreibung   & 8 Inhalt & Inhalt & Inhalt \tabularnewline
\bfseries Beschreibung   & 9 Inhalt & Inhalt & Inhalt \tabularnewline
\bfseries Beschreibung   & 0 Inhalt & Inhalt & Inhalt \tabularnewline
\bfseries Beschreibung   & 1 Inhalt & Inhalt & Inhalt \tabularnewline
\bfseries Beschreibung   & 2 Inhalt & Inhalt & Inhalt \tabularnewline
\bfseries Beschreibung   & 3 Inhalt & Inhalt & Inhalt \tabularnewline
\bfseries Beschreibung   & 4 Inhalt & Inhalt & Inhalt \tabularnewline
\bfseries Beschreibung   & 5 Inhalt & Inhalt & Inhalt \tabularnewline
\bfseries Beschreibung   & 6 Inhalt & Inhalt & Inhalt \tabularnewline
\bfseries Beschreibung   & 7 Inhalt & Inhalt & Inhalt \tabularnewline
\bfseries Beschreibung   & 8 Inhalt & Inhalt & Inhalt \tabularnewline
\bfseries Beschreibung   & 9 Inhalt & Inhalt & Inhalt \tabularnewline
\bfseries Beschreibung   & 0 Inhalt & Inhalt & Inhalt \tabularnewline
\bfseries Beschreibung   & 1 Inhalt & Inhalt & Inhalt \tabularnewline
\bfseries Beschreibung   & 2 Inhalt & Inhalt & Inhalt \tabularnewline
\bfseries Beschreibung   & 3 Inhalt & Inhalt & Inhalt \tabularnewline
\bfseries Beschreibung   & 4 Inhalt & Inhalt & Inhalt \tabularnewline
\bfseries Beschreibung   & 5 Inhalt & Inhalt & Inhalt \tabularnewline
\bfseries Beschreibung   & 6 Inhalt & Inhalt & Inhalt \tabularnewline
\bfseries Beschreibung   & 7 Inhalt & Inhalt & Inhalt \tabularnewline
\bfseries Beschreibung   & 8 Inhalt & Inhalt & Inhalt \tabularnewline
\bfseries Beschreibung   & 9 Inhalt & Inhalt & Inhalt \tabularnewline
\bfseries Beschreibung   & 0 Inhalt & Inhalt & Inhalt \tabularnewline
\end{longtable}